% --- CORRECCIÓN IMPORTANTE EN LA PRIMERA LÍNEA ---
% Agregamos "xcolor=table" para que funcione \rowcolor sin errores
\documentclass[aspectratio=169, 10pt, xcolor=table]{beamer}

% --- Configuración del Tema (estilo profesional como el ejemplo) ---
\usetheme{Madrid}
\usecolortheme{dolphin}

% --- Paquetes necesarios ---
\usepackage[utf8]{inputenc}
\usepackage[spanish]{babel}
\usepackage{amsmath, amssymb}
\usepackage{booktabs}
\usepackage{graphicx}
\usepackage{listings}
\usepackage{tikz}
\usetikzlibrary{arrows.meta, positioning, shapes.geometric}

% --- Ruta de Imágenes (crea la carpeta "imagenes" y coloca ahí tus archivos) ---
\graphicspath{{./imagenes/}}

% --- Estilo para código Python ---
\definecolor{codegreen}{rgb}{0,0.6,0}
\definecolor{codegray}{rgb}{0.5,0.5,0.5}
\definecolor{codepurple}{rgb}{0.58,0,0.82}
\definecolor{backcolour}{rgb}{0.95,0.95,0.92}
\lstset{
    language=Python,
    backgroundcolor=\color{backcolour},
    commentstyle=\color{codegreen},
    keywordstyle=\color{blue},
    numberstyle=\tiny\color{codegray},
    stringstyle=\color{codepurple},
    basicstyle=\ttfamily\scriptsize,
    breaklines=true,
    frame=single,
    showstringspaces=false
}

% --- Pie de página personalizado ---
\makeatletter

\setbeamertemplate{footline}{
  \leavevmode%
  \hbox{%
  \begin{beamercolorbox}[wd=.333333\paperwidth,ht=2.25ex,dp=1ex,center]{author in head/foot}%
    \usebeamerfont{author in head/foot}\insertshortauthor
  \end{beamercolorbox}%
  \begin{beamercolorbox}[wd=.333333\paperwidth,ht=2.25ex,dp=1ex,center]{title in head/foot}%
    \usebeamerfont{title in head/foot}Sesión 3
  \end{beamercolorbox}%
  \begin{beamercolorbox}[wd=.333333\paperwidth,ht=2.25ex,dp=1ex,right]{date in head/foot}%
    \usebeamerfont{date in head/foot}Machine Learning \hfill \insertframenumber/\inserttotalframenumber \hspace*{0.3cm}
  \end{beamercolorbox}}%
  \vskip0pt%
}

\makeatother

% --- Datos de la portada ---
\title[SESIÓN 3]{\textbf{Sesión 3}}
\subtitle{Regresión Lineal y Regularización}
\author[Prof. C. Sánchez]{Carlos César Sánchez Coronel}
\institute{\textbf{MACHINE LEARNING}}
\date{}

% Logo opcional (descomentar si tienes la imagen)
% \titlegraphic{\includegraphics[height=1.5cm]{logo.png}}

% --- Eliminar la barra de navegación (opcional) ---
\setbeamertemplate{navigation symbols}{}

% --- Índice al inicio de cada sección ---
\AtBeginSection[]{
  \begin{frame}
    \frametitle{Contenido}
    \tableofcontents[currentsection]
  \end{frame}
}

% ============================================================================
% INICIO DEL DOCUMENTO
% ============================================================================
\begin{document}

% --- Portada ---
\begin{frame}
    \titlepage
\end{frame}

% --- Índice general ---
\begin{frame}{Contenido}
    \tableofcontents
\end{frame}

% ============================================================================
% SECCIÓN 1: REGRESIÓN LINEAL SIMPLE Y MÚLTIPLE
% ============================================================================
\section{Regresión Lineal Simple y Múltiple}

\begin{frame}{Regresión Lineal Simple}
    \begin{block}{Formulación}
        $y = \beta_0 + \beta_1 x + \varepsilon$
        \begin{itemize}
            \item $\beta_0$: intersección
            \item $\beta_1$: pendiente
            \item $\varepsilon$: error
        \end{itemize}
    \end{block}
    \begin{exampleblock}{Estimación MCO}
        \[
        \hat{\beta}_1 = \frac{\sum (x_i - \bar{x})(y_i - \bar{y})}{\sum (x_i - \bar{x})^2}, \quad
        \hat{\beta}_0 = \bar{y} - \hat{\beta}_1\bar{x}
        \]
    \end{exampleblock}
    \begin{itemize}
        \item \textbf{Interpretación:} $\hat{\beta}_1$ es el cambio esperado en $y$ por unidad de $x$.
        \item \textbf{Caso de uso:} Relación entre años de experiencia y salario.
    \end{itemize}
    % Basado en Pregunta 1
    % IMAGEN: Diagrama de dispersión con recta de regresión
\end{frame}

\begin{frame}{Regresión Lineal Múltiple}
    \begin{block}{Formulación matricial}
        $\mathbf{y} = \mathbf{X}\boldsymbol{\beta} + \boldsymbol{\varepsilon}$
        \begin{itemize}
            \item $\mathbf{y}$: $n \times 1$, $\mathbf{X}$: $n \times (p+1)$, $\boldsymbol{\beta}$: $(p+1) \times 1$, $\boldsymbol{\varepsilon}$: $n \times 1$
        \end{itemize}
    \end{block}
    \begin{exampleblock}{Solución MCO}
        $\hat{\boldsymbol{\beta}} = (\mathbf{X}^T\mathbf{X})^{-1}\mathbf{X}^T\mathbf{y}$
        \begin{itemize}
            \item Requiere que $\mathbf{X}^T\mathbf{X}$ sea invertible (no multicolinealidad perfecta).
        \end{itemize}
    \end{exampleblock}
    \begin{itemize}
        \item \textbf{Interpretación ceteris paribus:} $\beta_j$ mide el efecto de $x_j$ manteniendo las demás constantes.
        \item \textbf{Caso de uso:} Predicción de precios de viviendas (área, habitaciones, antigüedad, ubicación).
    \end{itemize}
    % Basado en Pregunta 2
    % IMAGEN: Tabla de coeficientes
\end{frame}

% ============================================================================
% SECCIÓN 2: SUPUESTOS DEL MODELO
% ============================================================================
\section{Supuestos del Modelo}

\begin{frame}{Supuestos de la Regresión Lineal}
    \begin{columns}[t]
        \column{0.5\textwidth}
        \begin{alertblock}{Linealidad}
            Relación lineal entre predictores y respuesta. \\
            \textbf{Detección:} gráfico de residuos vs. ajustados.
        \end{alertblock}
        \begin{alertblock}{Independencia}
            Errores no correlacionados. \\
            \textbf{Detección:} gráfico de residuos vs. orden.
        \end{alertblock}
        \column{0.5\textwidth}
        \begin{alertblock}{Homocedasticidad}
            Varianza constante de los errores. \\
            \textbf{Detección:} patrón de embudo en residuos vs. ajustados.
        \end{alertblock}
        \begin{alertblock}{Normalidad}
            Errores ∼ normal (para inferencia). \\
            \textbf{Detección:} Q-Q plot, histograma.
        \end{alertblock}
    \end{columns}
    \vspace{0.3cm}
    \begin{itemize}
        \item Para predicción, la normalidad es menos crítica.
        \item La violación de homocedasticidad afecta intervalos de confianza.
    \end{itemize}
    % Basado en Pregunta 3
    % IMAGEN: Gráficos de diagnóstico
\end{frame}

% ============================================================================
% SECCIÓN 3: MÉTRICAS DE EVALUACIÓN
% ============================================================================
\section{Métricas de Evaluación}

\begin{frame}{Métricas de Error en Regresión}
    \begin{columns}[t]
        \column{0.5\textwidth}
        \begin{block}{MSE (Error Cuadrático Medio)}
            $\displaystyle \frac{1}{n}\sum (y_i - \hat{y}_i)^2$
        \end{block}
        \begin{block}{RMSE (Raíz del MSE)}
            $\displaystyle \sqrt{\frac{1}{n}\sum (y_i - \hat{y}_i)^2}$
        \end{block}
        \column{0.5\textwidth}
        \begin{block}{MAE (Error Absoluto Medio)}
            $\displaystyle \frac{1}{n}\sum |y_i - \hat{y}_i|$
        \end{block}
        \begin{block}{$R^2$ (Coeficiente de Determinación)}
            $\displaystyle 1 - \frac{RSS}{TSS}$
        \end{block}
    \end{columns}
    \vspace{0.3cm}
    \begin{exampleblock}{Ejemplo: errores = [10,10,10,100]}
        MAE = 32.5, \quad RMSE ≈ 50.7
        \begin{itemize}
            \item RMSE penaliza más los errores grandes.
        \end{itemize}
    \end{exampleblock}
    % Basado en Pregunta 4 y 5
    % IMAGEN: Comparación de métricas
\end{frame}

\begin{frame}{Cuándo usar cada métrica}
    \begin{itemize}
        \item \textbf{RMSE:} Cuando los errores grandes son críticos (ej. rotura de stock).
        \item \textbf{MAE:} Robustez a outliers, todos los errores igual de importantes.
        \item \textbf{MSE:} Útil como función de pérdida (diferenciable).
        \item \textbf{$R^2$:} Bondad de ajuste, pero no sirve para comparar modelos con diferente número de predictores (usar $R^2$ ajustado).
    \end{itemize}
    \vspace{0.3cm}
    \begin{block}{$R^2$ ajustado}
        $\displaystyle R^2_{\text{adj}} = 1 - \frac{RSS/(n-p-1)}{TSS/(n-1)}$ penaliza la inclusión de variables irrelevantes.
    \end{block}
    % Basado en Pregunta 17
\end{frame}

% ============================================================================
% SECCIÓN 4: SOBREAJUSTE Y REGULARIZACIÓN
% ============================================================================
\section{Sobreajuste y Regularización}

\begin{frame}{El problema del sobreajuste}
    \begin{columns}[t]
        \column{0.5\textwidth}
        \begin{block}{Alta varianza}
            Modelo demasiado complejo, se ajusta al ruido.
            \begin{itemize}
                \item Buen rendimiento en entrenamiento, malo en prueba.
            \end{itemize}
        \end{block}
        \column{0.5\textwidth}
        \begin{block}{Alto sesgo}
            Modelo demasiado simple, no captura la señal.
            \begin{itemize}
                \item Mal rendimiento en ambos conjuntos.
            \end{itemize}
        \end{block}
    \end{columns}
    \vspace{0.3cm}
    \begin{exampleblock}{Trade-off sesgo-varianza}
        La regularización aumenta ligeramente el sesgo para reducir drásticamente la varianza, mejorando la generalización.
    \end{exampleblock}
    % Basado en Pregunta 6
    % IMAGEN: Gráfico sesgo-varianza
\end{frame}

% ============================================================================
% SECCIÓN 5: RIDGE (L2)
% ============================================================================
\section{Ridge (L2)}

\begin{frame}{Regularización Ridge}
    \begin{block}{Función de coste}
        $\displaystyle J_{\text{ridge}}(\boldsymbol{\beta}) = \sum_{i=1}^n (y_i - \hat{y}_i)^2 + \lambda \sum_{j=1}^p \beta_j^2$
    \end{block}
    \begin{itemize}
        \item $\lambda \ge 0$ controla la fuerza de regularización.
        \item Solución analítica: $\hat{\boldsymbol{\beta}}_{\text{ridge}} = (\mathbf{X}^T\mathbf{X} + \lambda\mathbf{I})^{-1}\mathbf{X}^T\mathbf{y}$
        \item \textbf{Efecto:} encoge los coeficientes hacia cero, pero ninguno se anula exactamente.
    \end{itemize}
    \begin{exampleblock}{Caso de uso}
        Muchas variables correlacionadas (ej. marketing con variables demográficas y de consumo). Ridge estabiliza las estimaciones.
    \end{exampleblock}
    % Basado en Pregunta 7 y 13
    % IMAGEN: Curva de regularización (coeficientes vs λ)
\end{frame}

% ============================================================================
% SECCIÓN 6: LASSO (L1)
% ============================================================================
\section{Lasso (L1)}

\begin{frame}{Regularización Lasso}
    \begin{block}{Función de coste}
        $\displaystyle J_{\text{lasso}}(\boldsymbol{\beta}) = \sum_{i=1}^n (y_i - \hat{y}_i)^2 + \lambda \sum_{j=1}^p |\beta_j|$
    \end{block}
    \begin{itemize}
        \item La penalización L1 puede forzar algunos coeficientes exactamente a cero.
        \item \textbf{Selección automática de variables:} produce modelos dispersos e interpretables.
        \item No tiene solución analítica cerrada (se usa descenso por coordenadas).
    \end{itemize}
    \begin{exampleblock}{Caso de uso}
        Identificar factores clave en un proceso industrial con cientos de sensores.
    \end{exampleblock}
    % Basado en Pregunta 8 y 11
    % IMAGEN: Comparación de regiones de restricción L1 vs L2
\end{frame}

% ============================================================================
% SECCIÓN 7: ELASTIC NET
% ============================================================================
\section{Elastic Net}

\begin{frame}{Elastic Net: combinación de L1 y L2}
    \begin{block}{Función de coste}
        $\displaystyle J_{\text{elastic}}(\boldsymbol{\beta}) = \sum_{i=1}^n (y_i - \hat{y}_i)^2 + \lambda \left( \alpha \sum_{j=1}^p |\beta_j| + \frac{1-\alpha}{2} \sum_{j=1}^p \beta_j^2 \right)$
    \end{block}
    \begin{itemize}
        \item $\alpha \in [0,1]$ controla la mezcla: $\alpha=0$ → Ridge, $\alpha=1$ → Lasso.
        \item \textbf{Ventajas:} selecciona grupos de variables correlacionadas, maneja $p > n$.
    \end{itemize}
    \begin{exampleblock}{Caso de uso}
        Análisis de expresión génica (miles de genes, pocas muestras). Elastic Net identifica grupos de genes relacionados con una enfermedad.
    \end{exampleblock}
    % Basado en Pregunta 9 y 12
    % IMAGEN: Diagrama de Venn de los tres métodos
\end{frame}

% ============================================================================
% SECCIÓN 8: SELECCIÓN DE HIPERPARÁMETROS
% ============================================================================
\section{Selección de Hiperparámetros}

\begin{frame}{Validación cruzada para elegir $\lambda$ y $\alpha$}
    \begin{columns}[t]
        \column{0.6\textwidth}
        \begin{block}{Proceso}
            \begin{enumerate}
                \item Dividir datos en $k$ folds.
                \item Para cada combinación de $\lambda$ (y $\alpha$), entrenar en $k-1$ folds y evaluar en el fold restante.
                \item Promediar el error de validación.
                \item Elegir el valor que minimiza el error.
            \end{enumerate}
        \end{block}
        \column{0.4\textwidth}
        \begin{exampleblock}{Métrica recomendada}
            RMSE o MAE según el contexto.
        \end{exampleblock}
    \end{columns}
    \vspace{0.3cm}
    \begin{itemize}
        \item \textbf{Curva de regularización:} muestra cómo varían los coeficientes con $\lambda$.
    \end{itemize}
    % Basado en Pregunta 10
    % IMAGEN: Curva de validación (error vs λ)
\end{frame}

% ============================================================================
% SECCIÓN 9: COMPARACIÓN DE MÉTODOS
% ============================================================================
\section{Comparación de Métodos}

\begin{frame}{Comparación: Ridge vs Lasso vs Elastic Net}
    \begin{table}
        \centering
        \renewcommand{\arraystretch}{1.3}
        \begin{tabular}{l c c c}
            \toprule
            \textbf{Característica} & \textbf{Ridge (L2)} & \textbf{Lasso (L1)} & \textbf{Elastic Net} \\
            \midrule
            Selección de variables & No & Sí (individual) & Sí (grupos) \\
            Manejo de multicolinealidad & Excelente & Regular & Bueno \\
            Estabilidad & Alta & Puede ser inestable & Buena \\
            Interpretabilidad & Media & Alta & Alta \\
            \bottomrule
        \end{tabular}
    \end{table}
    \begin{itemize}
        \item \textbf{Uso típico:} Ridge para variables correlacionadas, Lasso para modelos parsimoniosos, Elastic Net como compromiso.
    \end{itemize}
    % Basado en Pregunta 12
\end{frame}

% ============================================================================
% SECCIÓN 10: CASO DE ESTUDIO INTEGRADO
% ============================================================================
\section{Caso de Estudio Integrado}

\begin{frame}{Predicción de precios de viviendas (Ames, Kaggle)}
    \begin{block}{Datos}
        80 variables predictoras (área, habitaciones, año, ubicación, calidad...).
    \end{block}
    \begin{columns}[t]
        \column{0.5\textwidth}
        \textbf{Preprocesamiento}
        \begin{itemize}
            \item Imputación de nulos (mediana, moda).
            \item Codificación one-hot (∼200 features).
            \item Estandarización (necesaria para regularización).
        \end{itemize}
        \column{0.5\textwidth}
        \textbf{Modelado}
        \begin{itemize}
            \item MCO (sin regularización)
            \item Ridge (validación cruzada)
            \item Lasso (validación cruzada)
            \item Elastic Net (búsqueda de $\lambda$, $\alpha$)
        \end{itemize}
    \end{columns}
    \vspace{0.3cm}
    \begin{exampleblock}{Resultados esperados}
        MCO sobreajusta; Ridge estabiliza; Lasso selecciona 20-30 variables clave; Elastic Net balancea.
    \end{exampleblock}
    % Basado en Pregunta 16
    % IMAGEN: Tabla de resultados (RMSE, MAE, R²)
\end{frame}

% ============================================================================
% SECCIÓN 11: CASO DE ESTUDIO AVANZADO (Pregunta 20)
% ============================================================================
\section{Caso Avanzado: Fintech Hipotecaria}

\begin{frame}{Diseño de un sistema de valoración automática}
    \begin{block}{Requisitos}
        \begin{itemize}
            \item 500,000 registros, 100 variables.
            \item Alta precisión e interpretabilidad.
            \item Variables correlacionadas.
        \end{itemize}
    \end{block}
    \begin{columns}[t]
        \column{0.5\textwidth}
        \textbf{Pipeline}
        \begin{itemize}
            \item Imputación robusta.
            \item Log/Box-Cox para asimetría.
            \item Target encoding para ubicación.
            \item Interacciones (área×calidad).
            \item Estandarización.
        \end{itemize}
        \column{0.5\textwidth}
        \textbf{Modelo elegido: Elastic Net}
        \begin{itemize}
            \item Selecciona variables (interpretabilidad).
            \item Maneja correlaciones.
            \item Estable.
        \end{itemize}
    \end{columns}
    \vspace{0.2cm}
    \begin{itemize}
        \item Validación cruzada con RMSE.
        \item Monitoreo en producción: PSI, drift.
        \item Automatización MLOps.
    \end{itemize}
    % Basado en Pregunta 20
\end{frame}

% ============================================================================
% SECCIÓN 12: RESUMEN Y CONEXIÓN
% ============================================================================
\section{Resumen y Conexión}

\begin{frame}{Lo que hemos aprendido}
    \begin{exampleblock}{Conceptos clave}
        \begin{itemize}
            \item Regresión lineal: formulación, estimación, interpretación.
            \item Supuestos y métricas de evaluación.
            \item Sobreajuste y trade-off sesgo-varianza.
            \item Regularización: Ridge (L2), Lasso (L1), Elastic Net.
            \item Selección de hiperparámetros con validación cruzada.
            \item Aplicación a casos reales (viviendas, fintech, bioinformática).
        \end{itemize}
    \end{exampleblock}
    \vspace{0.3cm}
    \begin{block}{Conexión con el resto del curso}
        Estos conceptos se extienden a:
        \begin{itemize}
            \item Regresión logística regularizada.
            \item Redes neuronales (weight decay).
            \item Selección de características en modelos complejos.
        \end{itemize}
    \end{block}
\end{frame}

% --- Diapositiva final ---
\begin{frame}
    \centering
    \Huge \textbf{¡Gracias!}

\end{frame}

\end{document}